\documentclass[12pt]{article}
% Préambule commun à tous les documents extraits de « La Revue CAPES »
\usepackage[T1]{fontenc}
\usepackage[utf8]{inputenc}
\usepackage{lmodern}
\usepackage{amsmath,amssymb,amsthm,amsfonts,mathrsfs}
\usepackage{setspace}
\usepackage{listings}
\usepackage{pgfplots}
\pgfplotsset{compat=1.18}
\usepackage{longtable}
\usepackage{adjustbox}
\usepackage{lettrine}
\usepackage{tkz-tab}
\usepackage{changepage}
\usepackage{pdflscape}
\usepackage{graphicx}
\usepackage{tikz}
\usepackage{xcolor}
\usepackage[a4paper,top=2cm,bottom=2.5cm,left=2.5cm,right=2cm]{geometry}
\usepackage{fancyhdr}
\usepackage{draftwatermark}
\SetWatermarkText{La RevueCapes}
\SetWatermarkLightness{0.9}
\SetWatermarkScale{2}
\usepackage{hyperref}
\hypersetup{colorlinks=true,linkcolor=blue!50!black,urlcolor=blue!50!black}

\definecolor{revue}{RGB}{20,60,120}

\pagestyle{fancy}
\fancyhf{}
\renewcommand{\headrulewidth}{0.4pt}
\setlength{\headheight}{15pt}

% \entete{Matière}{Titre}{Type de document}
\newcommand{\entete}[3]{%
  \fancyhead[L]{\small\textsc{La Revue CAPES} -- #1}%
  \fancyhead[R]{\small #2}%
  \fancyfoot[C]{\thepage}%
  \thispagestyle{plain}%
  \begin{center}
    {\color{revue}\rule{\textwidth}{1.2pt}}\\[4pt]
    {\small\textsc{Concours directs CAP-CEG / CAPES -- Burkina Faso}}\\[6pt]
    {\LARGE\bfseries\color{revue} #2}\\[6pt]
    {\large\itshape #1 -- #3}\\[2pt]
    {\color{revue}\rule{\textwidth}{1.2pt}}
  \end{center}
  \vspace{0.5cm}%
}

% Encadré pour les remarques ajoutées dans les corrigés
\newcommand{\remarque}[1]{\par\medskip\noindent\fcolorbox{revue}{revue!6}{\parbox{\dimexpr\linewidth-2\fboxsep-2\fboxrule}{\textbf{\color{revue}Remarque.} #1}}\par\medskip}


\begin{document}
\entete{Culture générale}{Corrigé du sujet type de dissertation n°4}{Correction}
\begin{spacing}{1.5}

\noindent\textbf{Sujet.} « La lutte contre le terrorisme au Burkina Faso doit être envisagée en amont avec les forces armées de défense et de sécurité, et en aval avec l'appui des populations. » Justifiez ce point de vue avec trois (03) arguments et donnez trois (03) limites.

\rubrique{1. Analyse du sujet}

\begin{itemize}
\item \textbf{Type de sujet :} la consigne impose la structure : une première partie avec \textbf{trois arguments} qui justifient la thèse, une seconde avec \textbf{trois limites}. Chaque argument et chaque limite doit être expliqué et illustré par un exemple.
\item \textbf{Mots-clés :}
  \begin{itemize}
  \item « en amont » : en première ligne, dans l'action directe (opérations militaires, sécurisation du territoire) ;
  \item « forces armées de défense et de sécurité » (FDS) : armée, gendarmerie, police ;
  \item « en aval » : en soutien, en accompagnement de l'action militaire ;
  \item « l'appui des populations » : renseignement, vigilance, solidarité, engagement civique (par exemple les Volontaires pour la défense de la patrie, VDP).
  \end{itemize}
\item \textbf{Problématique :} pourquoi la lutte antiterroriste doit-elle associer les FDS et les populations, et quelles sont les difficultés de cette collaboration ?
\end{itemize}

\rubrique{2. Plan détaillé}

\begin{enumerate}
\item[I.] \textbf{Trois arguments en faveur de cette approche.} 1. Les FDS ont la compétence et les moyens de la riposte armée. 2. Les populations fournissent le renseignement indispensable. 3. La cohésion entre FDS et populations prive les terroristes de leur base sociale.
\item[II.] \textbf{Trois limites.} 1. La méfiance entre FDS et certaines populations. 2. L'exposition des civils aux représailles et les risques de dérives. 3. L'insuffisance des moyens et des réponses non militaires.
\end{enumerate}

\rubrique{3. Proposition de rédaction}

\partie{Introduction}

Depuis 2015, le Burkina Faso fait face à des attaques terroristes récurrentes qui ont coûté la vie à des milliers de civils et de soldats et contraint plus de deux millions de personnes à quitter leurs foyers. Face à cette menace, une conviction s'est imposée : « la lutte contre le terrorisme au Burkina Faso doit être envisagée en amont avec les forces armées de défense et de sécurité, et en aval avec l'appui des populations ». Autrement dit, l'action militaire, conduite par les professionnels de la défense, doit être soutenue par l'engagement des citoyens. Pourquoi cette double mobilisation est-elle nécessaire, et quelles en sont les limites ? Nous justifierons d'abord ce point de vue par trois arguments, avant d'en présenter trois limites.

\partie{I. Trois arguments en faveur de cette approche}

\souspartie{Argument 1 : les FDS ont la compétence et les moyens de la riposte armée}
La lutte contre des groupes armés qui disposent d'armes de guerre, de motos et d'explosifs exige une force organisée, formée et équipée. Seules les FDS ont la mission constitutionnelle, la formation, la discipline et les moyens (renseignement, logistique, appui aérien) pour mener des opérations de reconquête, protéger les axes routiers, escorter les convois et sécuriser les villes. Ce sont elles qui doivent être « en amont », en première ligne. \emph{Exemple :} les opérations de sécurisation et d'escorte des convois de ravitaillement vers les localités sous blocus, comme Djibo.

\souspartie{Argument 2 : les populations fournissent le renseignement indispensable}
Les terroristes se fondent dans la population et utilisent la surprise. Les habitants connaissent leur territoire, remarquent les mouvements suspects, les étrangers, les caches. Sans leurs informations, les FDS combattent « à l'aveugle ». Le renseignement humain est l'arme principale d'une guerre asymétrique. \emph{Exemple :} les numéros verts mis à la disposition du public et les comités locaux de vigilance permettent de signaler à temps les menaces.

\souspartie{Argument 3 : la cohésion entre FDS et populations prive les terroristes de leur base sociale}
Les groupes armés cherchent à recruter, à se ravitailler et à se cacher parmi les populations. Lorsque celles-ci font bloc avec les FDS, ces groupes perdent leurs soutiens, leurs recrues et leurs ressources. L'appui des populations prend aussi la forme de l'engagement volontaire : les Volontaires pour la défense de la patrie (VDP), créés en 2020, combattent aux côtés des FDS ; la solidarité envers les personnes déplacées internes et l'effort de guerre (contributions au Fonds de soutien patriotique) renforcent la résilience de la nation. \emph{Exemple :} la mobilisation générale et la participation de nombreux jeunes au recrutement des VDP.

\transition{Cette double mobilisation est donc nécessaire. Elle comporte toutefois des limites qu'il faut regarder en face.}

\partie{II. Trois limites}

\souspartie{Limite 1 : la méfiance entre FDS et certaines populations}
Dans certaines zones, la confiance entre FDS et populations est fragile. Des allégations d'exactions contre des civils, des arrestations arbitraires ou la stigmatisation de certaines communautés poussent une partie de la population à se taire, voire à rejoindre les groupes armés. Sans confiance, il n'y a pas de renseignement. \emph{Exemple :} les accusations de violences commises contre des civils au nom de la lutte antiterroriste, régulièrement dénoncées par des organisations de défense des droits humains.

\souspartie{Limite 2 : l'exposition des civils aux représailles et les risques de dérives}
Les populations qui collaborent avec les FDS deviennent des cibles : les groupes terroristes punissent cruellement les villages accusés de coopération. Les VDP, moins formés et moins équipés que les soldats professionnels, paient un lourd tribut. Par ailleurs, l'armement de civils comporte des risques : règlements de comptes, dérives communautaires, difficulté de contrôle. \emph{Exemple :} des attaques ont visé des localités en représailles à la présence de VDP.

\souspartie{Limite 3 : l'insuffisance des moyens et des réponses non militaires}
La collaboration FDS-populations ne suffit pas si les moyens manquent (équipement, formation, prise en charge des blessés et des familles) et si les causes profondes du terrorisme ne sont pas traitées. Pauvreté, chômage des jeunes, absence de l'État, conflits fonciers, radicalisation : ces problèmes appellent aussi des réponses politiques, économiques, éducatives et judiciaires. Une approche uniquement sécuritaire risque de traiter les symptômes sans guérir le mal.

\partie{Conclusion}

La lutte contre le terrorisme au Burkina Faso doit effectivement associer, en amont, les forces armées de défense et de sécurité, seules capables de mener la riposte armée, et, en aval, les populations, dont le renseignement, l'engagement et la cohésion privent les terroristes de leurs soutiens. Cette approche connaît cependant des limites : la méfiance entre certaines populations et les FDS, l'exposition des civils aux représailles et aux dérives, et l'insuffisance d'une réponse uniquement sécuritaire. Pour être efficace, elle doit donc s'accompagner du respect des droits humains, de la protection des civils, et d'une politique de développement et de justice capable de tarir les sources du recrutement.

\end{spacing}
\end{document}
